\makeatletter
\def\input@path{{./}{../}{../../}{../../../}{../../../../}{preamble/}{../preamble/}{../../preamble/}{../../../preamble/}{../../../../preamble/}{book2026-2027/preamble/}}
\makeatother
\input{preamble_exams.tex}
\input{graphs/G11_Histograms.tex}
\graphicspath{{./}{../}{../../}{../../../}{../../../../}{preamble/}{../preamble/}{../../preamble/}{../../../preamble/}{../../../../preamble/}{book2026-2027/preamble/}}
\setYColumnWidth{0.27}
\setlength{\columnsep}{6mm}
\usepackage{needspace}
\NewCommandCopy\SharedSubsectionBox\SubsectionBox
% Preserve the shared numbered/starred forms within the local width adjustment.
\RenewDocumentCommand{\SubsectionBox}{s m}{%
  \par\needspace{5\baselineskip}%
  \begingroup\setlength{\linewidth}{\dimexpr\linewidth-2\fboxsep\relax}%
  \IfBooleanTF{#1}{\SharedSubsectionBox*{#2}}{\SharedSubsectionBox{#2}}\endgroup}
\hypersetup{pdftitle={Catch ups exam}}

\setlength{\fboxsep}{1pt}

\begin{document}
\noindent
\begin{tabularx}{\textwidth}{|C{2.8cm}|C{\dimexpr\textwidth-5.6cm-6\tabcolsep-4\arrayrulewidth\relax}|C{2.8cm}|}
\hline
\includegraphics[width=2.5cm]{logo.png} &
\CellCenter{\textbf{GLOBAL ECONOMICS}\par
\textbf{Grade 11}\par
\textbf{Catch ups exam}\par
\textbf{Continuous distributions and sample data}\par
\textbf{Teacher: Nicolás López Cuéllar}} &
\CellCenter{\textbf{1st TERM}\par\textbf{2026--2027}}\\\hline
\end{tabularx}

\LearningObjective
{Identify probability distributions and their supports; use geometric interval probabilities and exponential waiting times; calculate theoretical moments, empirical probabilities, and sample statistics; and interpret errors relative to theoretical models in economic and operational contexts.}
{	\footnotesize
	\textbf{C1.} Characterises discrete and continuous random variables, identifies suitable distributions and supports, and interprets density and cumulative probability in context.\par
	\textbf{C2.} Applies uniform and triangular models and constant or piecewise-linear densities to interval probabilities using geometric areas.\par
	\textbf{C3.} Uses circular geometry and arc proportions to calculate probabilities for uniform timing and directions around a circle.\par
	\textbf{C4.} Applies exponential models to constant-rate waiting times, infers event rates, and assesses model suitability.\par
	\textbf{C5.} Calculates and interprets the mean, variance, and standard deviation of continuous distributions using standard formulas.\par
	\textbf{C6.} Estimates interval probabilities using sample relative frequencies and compares them with theoretical probabilities.\par
	\textbf{C7.} Calculates sample means, variances, and standard deviations for different sample sizes and data representations and compares them with theory.
}

\textbf{Name:} \rule{10cm}{0.4pt}\hfill\textbf{Date:} \rule{3cm}{0.4pt}
\vspace{-4mm}

\begin{multicols}{2}

\SubsectionBox*{Criteria assessment}
This exam assesses C1--C7. Each question and task shows its points. \textbf{Passing requires a minimum of 5 cumulative points per criterion, including points earned in previous learning activities}. Show all procedures; answers without the required procedure receive no credit. Partial credit may be awarded for a correct method with arithmetic errors. Missing names or late submissions incur a 1-point deduction per criterion.

\begin{center}
	\small\setlength{\tabcolsep}{5pt}
	\begin{tabular}{|c|c|c|c|c||c|c|c|c|}
		\hline
		\textbf{\textnumero}&\textbf{C1}&\textbf{C2}&\textbf{C3}&\textbf{C4}&
		\textbf{\textnumero}&\textbf{C5}&\textbf{C6}&\textbf{C7}\\\hline
		1&6&--&--&--&8&2&--&--\\\hline
		2&--&6&--&--&9&2&--&--\\\hline
		3&--&--&2&--&10&2&--&--\\\hline
		4&--&--&3&--&11&--&2&--\\\hline
		5&--&--&--&3&12&--&2&--\\\hline
		6&--&--&--&3&13&--&2&--\\\hline
		7&--&--&--&2&14&--&--&2\\\hline
		&&&&&15&--&--&2\\\hline
		&&&&&16&--&--&2\\\hline
		\textbf{Tot.}&\textbf{6}&\textbf{6}&\textbf{6}&\textbf{7}&
		\textbf{Tot.}&\textbf{6}&\textbf{6}&\textbf{6}\\\hline
	\end{tabular}
\end{center}


\raggedcolumns
\SubsectionBox{Models in organisations \score{C1-\sml{6}}}


Colombian organisations model two random variables in each context below. \textbf{For each variable}: classify it as discrete or continuous, and identify its distribution.

\textbf{1. Cooperative repayments.} A cooperative records (a) whether one randomly selected member repays a monthly instalment on time, coded as 1 if timely and 0 otherwise; and (b) the number of timely repayments among 40 members. Each member repays on time independently with probability 0.92. \score{1}

\textbf{2. Warehouse quality checks.} A warehouse records (a) the number of invoices checked until the first incorrect invoice, including that invoice, with independent checks and a constant error probability of 0.08; and (b) a scale's measurement error between $-3$ and $3$ grams, with equal-length intervals having equal probabilities and no errors outside these limits. \score{1}

\textbf{3. Printing orders.} A printer records (a) the number of orders received in one hour, arriving independently at a constant rate of 2 per hour, one at a time, with multiple arrivals in a sufficiently short interval negligible; and (b) a poster order's completion time on $[10,30]$ minutes, whose density rises linearly from zero at 10 to a single peak at 18, then falls linearly to zero at 30, with no values outside this range. \score{1}

\textbf{4. Coffee production.} A roaster records (a) the fraction of a storage tank used at dispatch on $[0,1]$, whose density increases along one straight line from zero at 0 to its maximum at 1 and is zero elsewhere; and (b) a filling machine's measurement error, modelled by an exact symmetric bell curve on the real line, centred at 0 grams with standard deviation 4 grams. \score{1}

\textbf{5. Digital shop operations.} A digital shop records (a) the number of refund requests among 50 purchases, each independently resulting in a request with probability 0.04; and (b) the number of customer enquiries received over three days, arriving independently at a constant rate of 5 per day, one at a time, with multiple arrivals in a sufficiently short interval negligible. \score{1}

\textbf{6. Urban delivery services.} A delivery company in Bogotá studies two aspects of its daily operations. (a) Delivery times range from 15 to 45 minutes. Short deliveries are the most common, and the likelihood of observing longer delivery times decreases steadily until reaching zero at the upper limit. (b) The company also monitors the weight of its packages. Most packages weigh close to 5 kg, while lighter and heavier packages become progressively less common in a symmetric pattern. The standard deviation of package weights is 1.2 kg. \score{1}

\begin{minipage}{\linewidth}
	\SubsectionBox{Market-stall setup \score{C2-\sml{6}}}
	The setup time $S$ for a stall at a Bogotá craft market ranges from 5 to 25 minutes. Its density is zero at both endpoints, rises linearly to its single maximum at 13 minutes, then falls linearly to zero; it is zero outside $[5,25]$.

	Find $P(9\le S\le19)$. \score{3}

	Suppose the range changes to $[0,30]$, with  other properties unchanged. Recalculate $P(9\le S\le19)$. \score{3}

\end{minipage}\par

\begin{minipage}{\linewidth}
\SubsectionBox{A repeating inspection cycle \score{C3-\sml{3}}}
A cold-chain alarm occurs at a uniformly distributed position within a repeating 60-minute inspection cycle. The cycle starts at minute 0 and restarts immediately after minute 60.

\begin{itemize}
	\item Find the probability that the alarm occurs in the window from minute 52 through the restart to minute 11 of the next cycle. \score{1}
	\item Find the probability that the alarm occurs within 200 seconds before or after the start of a cycle. \score{2}
\end{itemize}

\end{minipage}\par
\begin{minipage}{\linewidth}
\SubsectionBox{Direction of a survey drone \score{C3-\sml{3}}}
A survey drone's initial heading is uniformly distributed around a full circle. A reference marker points north. ``Within an angle'' means that tolerance on \textbf{either side} of north, including wraparound.

Compare the probability of a heading within $\pi/9$ radians of north with that of a heading within $15^\circ$ of north. Calculate both probabilities and state which exceeds the other and by how much. \score{3}

\end{minipage}\par
\begin{minipage}{\linewidth}
\SubsectionBox{Charging-space availability \score{C4-\sml{3}}}
Charging spaces at a logistics centre become free according to a Poisson process. The opening rate differs between the four scenarios below but remains constant within each scenario. In each case, observe from a fixed time and let $T$ be the waiting time in minutes for the next opening.

With an average of 5 openings every 30 minutes, find the probability of waiting less than 4 minutes, $P(T<4)$. \score{0.5}

With a mean waiting time of 4 minutes between openings, find the probability of waiting more than 6 minutes, $P(T>6)$. \score{0.5}

With an average of 3 openings every 20 minutes, find the probability of waiting between 2 and 7 minutes, $P(2<T<7)$. \score{1}

With an average of 12 openings per hour, find the probability that the waiting time is outside the range from 3 to 8 minutes, $P(T<3\text{ or }T>8)$. \score{1}

\end{minipage}\par
\begin{minipage}{\linewidth}
\SubsectionBox{Electronic receipt rates \score{C4-\sml{2}}}
A digital shop issues electronic receipts according to a Poisson process with a constant rate during each shift. The morning and afternoon shifts have different rates. From a fixed observation time within each shift, let $T_m$ and $T_a$ be the waiting times in minutes for the next receipt in the morning and afternoon, respectively. Assume the relevant shift continues throughout each waiting time.

In the morning, $P(T_m>10)=0.40$. Infer the morning rate in receipts per hour. \score{1}

In the afternoon, $P(T_a<6)=0.80$. Infer the afternoon rate in receipts per hour. \score{1}

\end{minipage}\par
\begin{minipage}{\linewidth}
\SubsectionBox{Two payment channels \score{C4-\sml{2}}}
Completed payments on channels A and B follow independent constant-rate Poisson processes. At a fixed observation time, the probability that A's next payment occurs within 5 minutes is 0.75. The probability that B's next payment takes more than 5 minutes is 0.50. Count payments on both channels from that same observation time.

Infer both rates. After how many minutes does A's expected payment count exceed B's expected count by 3? Find the time for this difference in expectations. \score{2}

\end{minipage}\par
\begin{minipage}{\linewidth}
\SubsectionBox{Two labelling stations \score{C5-\sml{2}}}
At a food cooperative, let $A$ and $B$ be labelling times in minutes at two stations. Both have support $[3,15]$.

For A, the density increases along one straight line from zero at 3 to its highest value at 15. For B, equal-length time intervals within the support have equal probabilities. Both densities are zero outside $[3,15]$.

Identify A's distribution and calculate its mean, variance, and standard deviation. \score{1}

Identify B's distribution and calculate its mean, variance, and standard deviation. \score{1}

\end{minipage}\par
\begin{minipage}{\linewidth}
\SubsectionBox{Display production and daily revenue \score{C5-\sml{2}}}
A workshop models display production time $D$ on $[1,10]$ hours. Its density decreases along one straight line from its highest value at 1 to zero at 10 and is zero outside this support.

The workshop models daily revenue $R$, measured in \textbf{thousands of Colombian pesos (COP)}, by an exact symmetric bell curve on the real line. Its density is centred at 1800 and has inflection points at 1680 and 1920, all in thousands of COP.

Identify D's distribution and calculate its mean, variance, and standard deviation. \score{1}

Identify R's distribution and infer its parameters. Calculate the mean and standard deviation in \textbf{COP} and the variance in \textbf{COP$^2$}. Interpret the mean in context. \score{1}

\end{minipage}\par
\begin{minipage}{\linewidth}
\SubsectionBox{Requests and repair times \score{C5-\sml{2}}}
A repair cooperative receives requests according to a constant-rate Poisson process averaging 12 requests per hour. Let $T$ be the interarrival time in minutes.

The duration $S$ of a repair has support $[10,30]$ minutes. Its density rises linearly from zero at 10 to a single peak at 18, then falls linearly to zero at 30; it is zero outside this range.

Identify T's distribution and calculate its mean and standard deviation in minutes and its variance in minutes$^2$. \score{1}

Identify S's distribution and calculate its mean and standard deviation in minutes and its variance in minutes$^2$. \score{1}

\end{minipage}\par


\SubsectionBox{Parcel masses \score{C6-\sml{2}}}
A sample of 50 parcel masses $X$, in kilograms, is recorded in two non-overlapping parts. The first 20 masses are listed individually:
\[
\begin{array}{rrrrr}
1.2&1.5&2.0&2.2&2.4\\
2.5&2.7&3.0&3.1&3.4\\
1.8&2.0&2.3&2.6&2.9\\
3.0&3.2&3.5&3.8&4.0
\end{array}
\]
The remaining 30 masses are summarised in a frequency table:
\[
\begin{array}{c|rrrrrr}
x_i&2.0&2.4&2.6&2.8&3.0&3.2\\\hline
f_i&3&5&7&6&5&4
\end{array}
\]

\newpage 

The model gives $P(2\le X<3)=0.40$; include 2 and exclude 3. Combine both parts as one sample.

Find the empirical probability of this event. \score{1}

Calculate its absolute and percentage errors. \score{1}


\begin{minipage}{\linewidth}
\SubsectionBox{File transfers, grouped data \score{C6-\sml{2}}}
Transfer times $X$ in seconds are grouped below. Classes include the left endpoint and exclude the right endpoint.
\[
\begin{array}{c|rrrr}
\text{Class}&[0,4)&[4,8)&[8,12)&[12,16)\\\hline
f_i&3&9&5&3
\end{array}
\]
The model gives $P(4\le X<12)=0.75$.

Find the empirical probability of this event. \score{1}

Calculate its absolute and percentage errors. \score{1}

\end{minipage}\par
\begin{minipage}{\linewidth}
\SubsectionBox{Collection queues, histogram \score{C6-\sml{2}}}
The histogram is the only supplied sample record of collection waiting times $X$ in minutes. Equal-width classes include the left endpoint and exclude the right endpoint. Bar heights are frequencies; the horizontal axis gives minutes.
\HistogramGraph{(0,2),(2,4),(4,6),(6,8)}{4,10,8,2}
{0.94\linewidth}{4.4cm}{blue!25}{2}{2}
The model gives $P(2\le X<6)=0.70$.

Find the empirical probability of this event. \score{1}

Calculate its absolute and percentage errors. \score{1}

\end{minipage}\par
\begin{minipage}{\linewidth}

\SubsectionBox{Cold-room readings, mixed data \score{C7-\sml{2}}}
A sample of 12 cold-room readings, in degrees Celsius, is recorded in two non-overlapping parts. The first five readings are $6,8,9,11,12$. The remaining seven readings are summarised below:
\[
\begin{array}{c|rrr}x_i&6&9&12\\\hline f_i&2&3&2\end{array}
\]
Combine both parts as one sample. The theoretical mean is $9\,{}^\circ\mathrm C$ and variance is $5\,({}^\circ\mathrm C)^2$.

Calculate mean, variance and standard deviation. \score{1}

Calculate absolute and percentage errors for mean and variance. \score{1}

\end{minipage}\par
\begin{minipage}{\linewidth}
\SubsectionBox{Courier routes, grouped data \score{C7-\sml{2}}}
Route times in minutes are grouped below. Classes include the left endpoint and exclude the right endpoint. The theoretical mean is 14 minutes and variance is 12 minutes$^2$.
\[
\begin{array}{c|rrr}
\text{Class}&[8,12)&[12,16)&[16,20)\\\hline
f_i&2&3&1
\end{array}
\]

Calculate mean, variance and standard deviation. \score{1}

Calculate absolute and percentage errors for mean and variance. \score{1}

\end{minipage}\par
\begin{minipage}{\linewidth}
\SubsectionBox{Reusable-container masses \score{C7-\sml{2}}}
The histogram is the only supplied sample record of container masses in kilograms. Equal-width classes include the left endpoint and exclude the right endpoint. Bar heights are frequencies; the horizontal axis gives kilograms. The theoretical mean is 7 kg and variance is 4 kg$^2$.
\HistogramGraph{(4,6),(6,8),(8,10),(10,12)}{2,3,2,1}
{0.94\linewidth}{4.4cm}{orange!30}{2}{1}

Calculate mean, variance and standard deviation. \score{1}

Calculate absolute and percentage errors for mean and variance. \score{1}

\end{minipage}\par
\end{multicols}
\end{document}

